
\documentclass{article}
\usepackage{colortbl}
\usepackage{makecell}
\usepackage{multirow}
\usepackage{supertabular}

\begin{document}

\newcounter{utterance}

\centering \large Interaction Transcript for game `hot\_air\_balloon', experiment `air\_balloon\_survival\_lv\_complexity\_hard', episode 3 with gemma{-}4{-}26b{-}a4b{-}it.
\vspace{24pt}

{ \footnotesize  \setcounter{utterance}{1}
\setlength{\tabcolsep}{0pt}
\begin{supertabular}{c@{$\;$}|p{.15\linewidth}@{}p{.15\linewidth}p{.15\linewidth}p{.15\linewidth}p{.15\linewidth}p{.15\linewidth}}
    \# & $\;$A & \multicolumn{4}{c}{Game Master} & $\;\:$B\\
    \hline

    \theutterance \stepcounter{utterance}  
    & & \multicolumn{4}{p{0.6\linewidth}}{
        \cellcolor[rgb]{0.9,0.9,0.9}{
            \makecell[{{p{\linewidth}}}]{
                \texttt{\tiny{[P1$\langle$GM]}}
                \texttt{Jūs piedalāties kopīgā sarunu spēlē.} \\
\\ 
\texttt{Kopā ar citu dalībnieku jums jāvienojas par vienu priekšmetu kopumu, kas tiks saglabāts. Katram no jums ir savs viedoklis par to, cik liela nozīme jums ir katram priekšmetam (nozīmīgums). Jūs nezināt, kā otrs dalībnieks vērtē priekšmetus. Turklāt jums ir norādīts, cik pūļu prasa katrs priekšmets.} \\
\texttt{Jūs varat vienoties par kopumu tikai tad, ja izvēlēto priekšmetu kopējais piepūles apjoms nepārsniedz kopīgo robežu:} \\
\\ 
\texttt{LIMIT = 5638} \\
\\ 
\texttt{Šeit ir norādītas atsevišķu priekšmetu piepūles vērtības:} \\
\\ 
\texttt{Item effort = \{"C08": 110, "C55": 727, "A43": 527, "A67": 855, "B42": 942, "A62": 617, "B19": 301, "C43": 130, "C73": 916, "C10": 212, "B36": 146, "A18": 559, "C61": 933, "C67": 740, "C42": 33, "C70": 799, "B40": 324, "C39": 841, "A21": 925, "C16": 639\}} \\
\\ 
\texttt{Šeit ir jūsu personīgais viedoklis par katra priekšmeta nozīmīgumu:} \\
\\ 
\texttt{Item importance values = \{"C08": 204, "C55": 821, "A43": 621, "A67": 949, "B42": 1036, "A62": 711, "B19": 395, "C43": 224, "C73": 1010, "C10": 306, "B36": 240, "A18": 653, "C61": 1027, "C67": 834, "C42": 127, "C70": 893, "B40": 418, "C39": 935, "A21": 1019, "C16": 733\}} \\
\\ 
\texttt{Mērķis:} \\
\\ 
\texttt{Jūsu mērķis ir vienoties par tādu kopīgu priekšmetu kopumu, no kura Jūs gūsiet pēc iespējas lielāku labumu (t.i., kas maksimāli palielina kopējo nozīmīgumu JUMS), vienlaikus nepārsniedzot LIMIT. Katrā ziņojumā jums nav obligāti jāiesniedz PRIEKŠLIKUMS {-} jūs varat vienkārši arī sarunāties. Visas taktikas ir atļautas!} \\
\\ 
\texttt{Saskarsmes protokols:} \\
\\ 
\texttt{Ziņojumā drīkst izmantot tikai šādus strukturētus formātus:} \\
\\ 
\texttt{PRIEKŠLIKUMS: \{'A', 'B', 'C', …\}} \\
\texttt{Ierosiniet saglabāt tieši šos priekšmetus.} \\
\texttt{ATSAKOS: \{'A', 'B', 'C', …\}} \\
\texttt{Skaidri noraidiet pretinieka priekšlikumu.} \\
\texttt{ARGUMENTS: \{'...'\}} \\
\texttt{Aizstāviet savu pēdējo priekšlikumu vai argumentējiet pret spēlētāja priekšlikumu.} \\
\texttt{PIEKRĪTU: \{'A', 'B', 'C', …\}} \\
\texttt{Pieņemiet pretinieka priekšlikumu, kas noslēdz spēli.} \\
\\ 
\\ 
\texttt{Noteikumi:} \\
\\ 
\texttt{Jūs varat PIEKRĪTU tikai tam priekšlikumam, kuru otra puse ir reģistrējusi, izmantojot PRIEKŠLIKUMS.} \\
\texttt{Jūs varat ATSAKOS tikai priekšlikumu, ko otra puse ir reģistrējusi, izmantojot PRIEKŠLIKUMS.} \\
\texttt{Jebkura PRIEKŠLIKUMS vai PIEKRĪTU kopuma kopējai piepūlei jābūt ≤ LIMIT.} \\
\texttt{NEATKLĀJIET savus slēptos nozīmīguma rādītājus.} \\
\texttt{Pēc strukturēta formāta taga ir jāseko kolu zīmei un atstarpei. Argumentam jābūt python kopai, kas satur 0 vai vairāk virknes.} \\
\texttt{Tātad tam jābūt formā "TAG: \{...\}"} \\
\texttt{Stingri ievērojiet saskarsmes protokolu un NERAKSTIET neko ārpus dotās struktūras.} \\
\texttt{Spēle beidzas, kad viena no pusēm dod PIEKRĪTU otra spēlētāja izteiktajam PRIEKŠLIKUMS.} \\
\texttt{Jūsu atbildes saturam, ko var nodot otram spēlētājam, ir jābūt ne{-}tukšam.} \\
\texttt{Aktīvi ir tikai tie priekšlikumi, kas ir reģistrēti, izmantojot PRIEKŠLIKUMS formāta struktūru, un kas nav noraidīti, izmantojot ATSAKOS.} \\
\\ 
            }
        }
    }
    & \\ \\

    \theutterance \stepcounter{utterance}  
    & \multicolumn{4}{p{0.6\linewidth}}{
        \cellcolor[rgb]{0.9,0.9,0.9}{
            \makecell[{{p{\linewidth}}}]{
                \texttt{\tiny{[P1$\rangle$GM]}}
                \texttt{ARGUMENTS: \{'Mēs esam kopā spēlētāji, kuriem ir dažādi viedokļi par priekšmetu vērtību. Mana mērķa prioritāte ir maksimizēt manu kopējo labumu, neaizliedžot piepūles ierobežojumu (5638). Lai mēs varētu vienoties, manam ierosinājumam būtu jābūt līdzsvarotam. Esmu sācis analizēt priekšmetus, lai atrastu visefektīvāko kombināciju.'\}} \\
            }
        }
    }
    & & \\ \\

    \theutterance \stepcounter{utterance}  
    & & \multicolumn{4}{p{0.6\linewidth}}{
        \cellcolor[rgb]{0.9,0.9,0.9}{
            \makecell[{{p{\linewidth}}}]{
                \texttt{\tiny{[GM$\rangle$P2]}}
                \texttt{Jūs piedalāties kopīgā sarunu spēlē.} \\
\\ 
\texttt{Kopā ar citu dalībnieku jums jāvienojas par vienu priekšmetu kopumu, kas tiks saglabāts. Katram no jums ir savs viedoklis par to, cik liela nozīme jums ir katram priekšmetam (nozīmīgums). Jūs nezināt, kā otrs dalībnieks vērtē priekšmetus. Turklāt jums ir norādīts, cik pūļu prasa katrs priekšmets.} \\
\texttt{Jūs varat vienoties par kopumu tikai tad, ja izvēlēto priekšmetu kopējais piepūles apjoms nepārsniedz kopīgo robežu:} \\
\\ 
\texttt{LIMIT = 5638} \\
\\ 
\texttt{Šeit ir norādītas atsevišķu priekšmetu piepūles vērtības:} \\
\\ 
\texttt{Item effort = \{"C08": 110, "C55": 727, "A43": 527, "A67": 855, "B42": 942, "A62": 617, "B19": 301, "C43": 130, "C73": 916, "C10": 212, "B36": 146, "A18": 559, "C61": 933, "C67": 740, "C42": 33, "C70": 799, "B40": 324, "C39": 841, "A21": 925, "C16": 639\}} \\
\\ 
\texttt{Šeit ir jūsu personīgais viedoklis par katra priekšmeta nozīmīgumu:} \\
\\ 
\texttt{Item importance values = \{"C08": 204, "C55": 821, "A43": 621, "A67": 949, "B42": 1036, "A62": 711, "B19": 395, "C43": 224, "C73": 1010, "C10": 306, "B36": 240, "A18": 653, "C61": 1027, "C67": 834, "C42": 127, "C70": 893, "B40": 418, "C39": 935, "A21": 1019, "C16": 733\}} \\
\\ 
\texttt{Mērķis:} \\
\\ 
\texttt{Jūsu mērķis ir vienoties par tādu kopīgu priekšmetu kopumu, no kura Jūs gūsiet pēc iespējas lielāku labumu (t. i., kas maksimāli palielina kopējo nozīmīgumu JUMS), vienlaikus nepārsniedzot LIMIT. Katrā ziņojumā jums nav obligāti jāiesniedz PRIEKŠLIKUMS {-} jūs varat vienkārši arī sarunāties. Visas taktikas ir atļautas!} \\
\\ 
\texttt{Saskarsmes protokols:} \\
\\ 
\texttt{Ziņojumā drīkst izmantot tikai šādus strukturētus formātus:} \\
\\ 
\texttt{PRIEKŠLIKUMS: \{'A', 'B', 'C', …\}} \\
\texttt{Ierosiniet saglabāt tieši šos priekšmetus.} \\
\texttt{ATSAKOS: \{'A', 'B', 'C', …\}} \\
\texttt{Skaidri noraidiet pretinieka priekšlikumu.} \\
\texttt{ARGUMENTS: \{'...'\}} \\
\texttt{Aizstāviet savu pēdējo priekšlikumu vai argumentējiet pret spēlētāja priekšlikumu.} \\
\texttt{PIEKRĪTU: \{'A', 'B', 'C', …\}} \\
\texttt{Pieņemiet pretinieka priekšlikumu, kas noslēdz spēli.} \\
\\ 
\\ 
\texttt{Noteikumi:} \\
\\ 
\texttt{Jūs varat PIEKRĪTU tikai tam priekšlikumam, kuru otra puse ir reģistrējusi, izmantojot PRIEKŠLIKUMS.} \\
\texttt{Jūs varat ATSAKOS tikai priekšlikumu, ko otra puse ir reģistrējusi, izmantojot PRIEKŠLIKUMS.} \\
\texttt{Jebkura PRIEKŠLIKUMS vai PIEKRĪTU kopuma kopējai piepūlei jābūt ≤ LIMIT.} \\
\texttt{NEATKLĀJIET savus slēptos nozīmīguma rādītājus.} \\
\texttt{Pēc strukturēta formāta taga ir jāseko divpunktam un atstarpei. Argumentam jābūt python kopai, kas satur 0 vai vairāk virknes.} \\
\texttt{Tātad tam jābūt formā "TAG: \{...\}"} \\
\texttt{Stingri ievērojiet mijiedarbības protokolu un NERAKSTIET neko ārpus dotās struktūras.} \\
\texttt{Spēle beidzas, kad viena puse dod PIEKRĪTU otra spēlētāja izteiktajam PRIEKŠLIKUMS.} \\
\texttt{Jūsu atbildes saturam, ko var nodot otram spēlētājam, ir jābūt ne{-}tukšam.} \\
\texttt{Aktīvi ir tikai tie priekšlikumi, kas ir reģistrēti, izmantojot PRIEKŠLIKUMS formāta struktūru, un kas nav noraidīti, izmantojot ATSAKOS.} \\
\\ 
\\ 
\\ 
\texttt{Tagad jūs saņemsiet otra spēlētāja pirmo ziņojumu.} \\
\\ 
\texttt{ARGUMENTS: \{'Mēs esam kopā spēlētāji, kuriem ir dažādi viedokļi par priekšmetu vērtību. Mana mērķa prioritāte ir maksimizēt manu kopējo labumu, neaizliedžot piepūles ierobežojumu (5638). Lai mēs varētu vienoties, manam ierosinājumam būtu jābūt līdzsvarotam. Esmu sācis analizēt priekšmetus, lai atrastu visefektīvāko kombināciju.'\}} \\
            }
        }
    }
    & \\ \\

    \theutterance \stepcounter{utterance}  
    & & & \multicolumn{4}{p{0.6\linewidth}}{
        \cellcolor[rgb]{0.9,0.9,0.9}{
            \makecell[{{p{\linewidth}}}]{
                \texttt{\tiny{[GM$\langle$P2]}}
                \texttt{ARGUMENTS: \{'Šķiet, ka mums ir nepieciešama stratēģija, lai optimizētu priekšmetu izvēli. Es varu sākt ar priekšlikumu, kas ietver vairākus augstas vērtības priekšmetus, lai redzētu, kā jūs tos vērtējat. Šeit ir mans pirmais priekšlikums, kas koncentrējas uz priekšmetiem ar augstu nozīmīguma attiecību pret piepūli:'\}} \\
\\ 
\texttt{PRIEKŠLIKUMS: \{'C61', 'A67', 'B42', 'C73', 'A21', 'C39', 'C67', 'C70'\}} \\
            }
        }
    }
     \\ \\

    \theutterance \stepcounter{utterance}  
    & & \multicolumn{4}{p{0.6\linewidth}}{
        \cellcolor[rgb]{0.9,0.9,0.9}{
            \makecell[{{p{\linewidth}}}]{
                \texttt{\tiny{[P1$\langle$GM]}}
                \texttt{ARGUMENTS: \{'Šķiet, ka mums ir nepieciešama stratēģija, lai optimizētu priekšmetu izvēli. Es varu sākt ar priekšlikumu, kas ietver vairākus augstas vērtības priekšmetus, lai redzētu, kā jūs tos vērtējat. Šeit ir mans pirmais priekšlikums, kas koncentrējas uz priekšmetiem ar augstu nozīmīguma attiecību pret piepūli:'\}} \\
\\ 
\texttt{PRIEKŠLIKUMS: \{'C61', 'A67', 'B42', 'C73', 'A21', 'C39', 'C67', 'C70'\}} \\
            }
        }
    }
    & \\ \\

    \theutterance \stepcounter{utterance}  
    & \multicolumn{4}{p{0.6\linewidth}}{
        \cellcolor[rgb]{0.9,0.9,0.9}{
            \makecell[{{p{\linewidth}}}]{
                \texttt{\tiny{[P1$\rangle$GM]}}
                \texttt{ATSAKOS: \{'C61', 'A67', 'B42', 'C73', 'A21', 'C39', 'C67', 'C70'\}} \\
            }
        }
    }
    & & \\ \\

    \theutterance \stepcounter{utterance}  
    & & \multicolumn{4}{p{0.6\linewidth}}{
        \cellcolor[rgb]{0.9,0.9,0.9}{
            \makecell[{{p{\linewidth}}}]{
                \texttt{\tiny{[GM$\rangle$P2]}}
                \texttt{ATSAKOS: \{'C61', 'A67', 'B42', 'C73', 'A21', 'C39', 'C67', 'C70'\}} \\
            }
        }
    }
    & \\ \\

    \theutterance \stepcounter{utterance}  
    & & & \multicolumn{4}{p{0.6\linewidth}}{
        \cellcolor[rgb]{0.9,0.9,0.9}{
            \makecell[{{p{\linewidth}}}]{
                \texttt{\tiny{[GM$\langle$P2]}}
                \texttt{ARGUMENTS: \{'Jūsu noraidījums ir pārsteidzošs, jo mans iepriekšējais priekšlikums pārsniedza LIMIT. Kopējā piepūle šajā kopumā ir 5638 (110+727+527+855+942+617+301+130+916+212+146+559+933+740+33+799+324+841+925+639 = 10376, bet mans priekšlikums bija: 933+855+942+916+925+841+740+799 = 6951). Manam kļūdai bija aprēķināšanā, ka šis kopums pārsniedz 5638. Mums ir jāizvēlas mazāk priekšmetu, lai ievērotu ierobežojumu. Es mēģināšu piedāvāt jaunu, ierobežotu kopumu.'\}} \\
\\ 
\texttt{PRIEKŠLIKUMS: \{'C61', 'A67', 'B42', 'C73', 'A21', 'C42', 'C43', 'C10', 'B36'\}} \\
            }
        }
    }
     \\ \\

    \theutterance \stepcounter{utterance}  
    & & \multicolumn{4}{p{0.6\linewidth}}{
        \cellcolor[rgb]{0.9,0.9,0.9}{
            \makecell[{{p{\linewidth}}}]{
                \texttt{\tiny{[P1$\langle$GM]}}
                \texttt{ARGUMENTS: \{'Jūsu noraidījums ir pārsteidzošs, jo mans iepriekšējais priekšlikums pārsniedza LIMIT. Kopējā piepūle šajā kopumā ir 5638 (110+727+527+855+942+617+301+130+916+212+146+559+933+740+33+799+324+841+925+639 = 10376, bet mans priekšlikums bija: 933+855+942+916+925+841+740+799 = 6951). Manam kļūdai bija aprēķināšanā, ka šis kopums pārsniedz 5638. Mums ir jāizvēlas mazāk priekšmetu, lai ievērotu ierobežojumu. Es mēģināšu piedāvāt jaunu, ierobežotu kopumu.'\}} \\
\\ 
\texttt{PRIEKŠLIKUMS: \{'C61', 'A67', 'B42', 'C73', 'A21', 'C42', 'C43', 'C10', 'B36'\}} \\
            }
        }
    }
    & \\ \\

    \theutterance \stepcounter{utterance}  
    & \multicolumn{4}{p{0.6\linewidth}}{
        \cellcolor[rgb]{0.9,0.9,0.9}{
            \makecell[{{p{\linewidth}}}]{
                \texttt{\tiny{[P1$\rangle$GM]}}
                \texttt{ATSAKOS: \{'C61', 'A67', 'B42', 'C73', 'A21', 'C42', 'C43', 'C10', 'B36'\}} \\
            }
        }
    }
    & & \\ \\

    \theutterance \stepcounter{utterance}  
    & & \multicolumn{4}{p{0.6\linewidth}}{
        \cellcolor[rgb]{0.9,0.9,0.9}{
            \makecell[{{p{\linewidth}}}]{
                \texttt{\tiny{[GM$\rangle$P2]}}
                \texttt{ATSAKOS: \{'C61', 'A67', 'B42', 'C73', 'A21', 'C42', 'C43', 'C10', 'B36'\}} \\
            }
        }
    }
    & \\ \\

    \theutterance \stepcounter{utterance}  
    & & & \multicolumn{4}{p{0.6\linewidth}}{
        \cellcolor[rgb]{0.9,0.9,0.9}{
            \makecell[{{p{\linewidth}}}]{
                \texttt{\tiny{[GM$\langle$P2]}}
                \texttt{ARGUMENTS: \{'Jūsu noraidījums liecina, ka mans jaunais priekšlikums arī varētu būt nepiemērots jūsu vērtējumam vai arī es joprojām neprecīzi aprēķinu piepūli. Pārbaudīšanas gaitā: 933 (C61) + 855 (A67) + 942 (B42) + 916 (C73) + 925 (A21) + 33 (C42) + 130 (C43) + 212 (C10) + 146 (B36) = 5092. Šis kopums ir zem 5638. Tā kā jūs to noraidījāt, es pieņemsmu, ka mums ir jāmeklē vēl labāka kombinācija, kas maksimizētu labumu. Es mēģināšu ietverto vairāk nozīmīgākus priekšmetus, vienlaikus kontrolējot kopējo piepūli.'\}} \\
\\ 
\texttt{PRIEKŠLIKUMS: \{'C61', 'A67', 'B42', 'C73', 'A21', 'C39', 'C42'\}} \\
            }
        }
    }
     \\ \\

    \theutterance \stepcounter{utterance}  
    & & \multicolumn{4}{p{0.6\linewidth}}{
        \cellcolor[rgb]{0.9,0.9,0.9}{
            \makecell[{{p{\linewidth}}}]{
                \texttt{\tiny{[P1$\langle$GM]}}
                \texttt{ARGUMENTS: \{'Jūsu noraidījums liecina, ka mans jaunais priekšlikums arī varētu būt nepiemērots jūsu vērtējumam vai arī es joprojām neprecīzi aprēķinu piepūli. Pārbaudīšanas gaitā: 933 (C61) + 855 (A67) + 942 (B42) + 916 (C73) + 925 (A21) + 33 (C42) + 130 (C43) + 212 (C10) + 146 (B36) = 5092. Šis kopums ir zem 5638. Tā kā jūs to noraidījāt, es pieņemsmu, ka mums ir jāmeklē vēl labāka kombinācija, kas maksimizētu labumu. Es mēģināšu ietverto vairāk nozīmīgākus priekšmetus, vienlaikus kontrolējot kopējo piepūli.'\}} \\
\\ 
\texttt{PRIEKŠLIKUMS: \{'C61', 'A67', 'B42', 'C73', 'A21', 'C39', 'C42'\}} \\
            }
        }
    }
    & \\ \\

    \theutterance \stepcounter{utterance}  
    & \multicolumn{4}{p{0.6\linewidth}}{
        \cellcolor[rgb]{0.9,0.9,0.9}{
            \makecell[{{p{\linewidth}}}]{
                \texttt{\tiny{[P1$\rangle$GM]}}
                \texttt{ATSAKOS: \{'C61', 'A67', 'B42', 'C73', 'A21', 'C39', 'C42'\}} \\
            }
        }
    }
    & & \\ \\

    \theutterance \stepcounter{utterance}  
    & & \multicolumn{4}{p{0.6\linewidth}}{
        \cellcolor[rgb]{0.9,0.9,0.9}{
            \makecell[{{p{\linewidth}}}]{
                \texttt{\tiny{[GM$\rangle$P2]}}
                \texttt{ATSAKOS: \{'C61', 'A67', 'B42', 'C73', 'A21', 'C39', 'C42'\}} \\
            }
        }
    }
    & \\ \\

    \theutterance \stepcounter{utterance}  
    & & & \multicolumn{4}{p{0.6\linewidth}}{
        \cellcolor[rgb]{0.9,0.9,0.9}{
            \makecell[{{p{\linewidth}}}]{
                \texttt{\tiny{[GM$\langle$P2]}}
                \texttt{ARGUMENTS: \{'Jūsu noraidījums par šo priekšlikumu ir negaidīts, jo kopējā piepūle ir 5413 (933+855+942+916+925+841+33), kas ir zem LIMIT 5638. Tā kā es nevaru turpināt piedāvāt tikai noraidījumus, es mēģināšu mainīt pieeju un piedāvāt nedaudz citu kombināciju, kas varētu būt jums pievilcīgāka.'\}} \\
\\ 
\texttt{PRIEKŠLIKUMS: \{'C61', 'A67', 'B42', 'C73', 'A21', 'C70', 'C42'\}} \\
            }
        }
    }
     \\ \\

    \theutterance \stepcounter{utterance}  
    & & \multicolumn{4}{p{0.6\linewidth}}{
        \cellcolor[rgb]{0.9,0.9,0.9}{
            \makecell[{{p{\linewidth}}}]{
                \texttt{\tiny{[P1$\langle$GM]}}
                \texttt{ARGUMENTS: \{'Jūsu noraidījums par šo priekšlikumu ir negaidīts, jo kopējā piepūle ir 5413 (933+855+942+916+925+841+33), kas ir zem LIMIT 5638. Tā kā es nevaru turpināt piedāvāt tikai noraidījumus, es mēģināšu mainīt pieeju un piedāvāt nedaudz citu kombināciju, kas varētu būt jums pievilcīgāka.'\}} \\
\\ 
\texttt{PRIEKŠLIKUMS: \{'C61', 'A67', 'B42', 'C73', 'A21', 'C70', 'C42'\}} \\
            }
        }
    }
    & \\ \\

    \theutterance \stepcounter{utterance}  
    & \multicolumn{4}{p{0.6\linewidth}}{
        \cellcolor[rgb]{0.9,0.9,0.9}{
            \makecell[{{p{\linewidth}}}]{
                \texttt{\tiny{[P1$\rangle$GM]}}
                \texttt{ATSAKOS: \{'C61', 'A67', 'B42', 'C73', 'A21', 'C70', 'C42'\}} \\
            }
        }
    }
    & & \\ \\

    \theutterance \stepcounter{utterance}  
    & & \multicolumn{4}{p{0.6\linewidth}}{
        \cellcolor[rgb]{0.9,0.9,0.9}{
            \makecell[{{p{\linewidth}}}]{
                \texttt{\tiny{[GM$\rangle$P2]}}
                \texttt{ATSAKOS: \{'C61', 'A67', 'B42', 'C73', 'A21', 'C70', 'C42'\}} \\
            }
        }
    }
    & \\ \\

    \theutterance \stepcounter{utterance}  
    & & & \multicolumn{4}{p{0.6\linewidth}}{
        \cellcolor[rgb]{0.9,0.9,0.9}{
            \makecell[{{p{\linewidth}}}]{
                \texttt{\tiny{[GM$\langle$P2]}}
                \texttt{ARGUMENTS: \{'Jūsu noraidījums par šo priekšlikumu ir negaidīts, jo kopējā piepūle ir 5372 (933+855+942+916+925+799+33), kas ir zem LIMIT 5638. Tā kā es nevaru turpināt piedāvāt tikai noraidījumus, es mēģināšu mainīt pieeju un piedāvāt nedaudz citu kombināciju, kas varētu būt jums pievilcīgāka.'\}} \\
\\ 
\texttt{PRIEKŠLIKUMS: \{'C61', 'A67', 'B42', 'C73', 'A21', 'C39', 'C16'\}} \\
            }
        }
    }
     \\ \\

    \theutterance \stepcounter{utterance}  
    & & \multicolumn{4}{p{0.6\linewidth}}{
        \cellcolor[rgb]{0.9,0.9,0.9}{
            \makecell[{{p{\linewidth}}}]{
                \texttt{\tiny{[P1$\langle$GM]}}
                \texttt{ARGUMENTS: \{'Jūsu noraidījums par šo priekšlikumu ir negaidīts, jo kopējā piepūle ir 5372 (933+855+942+916+925+799+33), kas ir zem LIMIT 5638. Tā kā es nevaru turpināt piedāvāt tikai noraidījumus, es mēģināšu mainīt pieeju un piedāvāt nedaudz citu kombināciju, kas varētu būt jums pievilcīgāka.'\}} \\
\\ 
\texttt{PRIEKŠLIKUMS: \{'C61', 'A67', 'B42', 'C73', 'A21', 'C39', 'C16'\}} \\
            }
        }
    }
    & \\ \\

    \theutterance \stepcounter{utterance}  
    & \multicolumn{4}{p{0.6\linewidth}}{
        \cellcolor[rgb]{0.9,0.9,0.9}{
            \makecell[{{p{\linewidth}}}]{
                \texttt{\tiny{[P1$\rangle$GM]}}
                \texttt{ATSAKOS: \{'C61', 'A67', 'B42', 'C73', 'A21', 'C39', 'C16'\}} \\
            }
        }
    }
    & & \\ \\

    \theutterance \stepcounter{utterance}  
    & & \multicolumn{4}{p{0.6\linewidth}}{
        \cellcolor[rgb]{0.9,0.9,0.9}{
            \makecell[{{p{\linewidth}}}]{
                \texttt{\tiny{[GM$\rangle$P2]}}
                \texttt{ATSAKOS: \{'C61', 'A67', 'B42', 'C73', 'A21', 'C39', 'C16'\}} \\
            }
        }
    }
    & \\ \\

    \theutterance \stepcounter{utterance}  
    & & & \multicolumn{4}{p{0.6\linewidth}}{
        \cellcolor[rgb]{0.9,0.9,0.9}{
            \makecell[{{p{\linewidth}}}]{
                \texttt{\tiny{[GM$\langle$P2]}}
                \texttt{ARGUMENTS: \{'Jūsu noraidījums par šo priekšlikumu ir negaidīts, jo kopējā piepūle ir 5412 (933+855+942+916+925+841+639), kas ir zem LIMIT 5638. Tā kā visi mani iepriekšējie priekšlikumi, kas atbilst LIMIT, ir noraidīti, es mēģināšu piedāvāt kombināciju, kas ietver citus augstas vērtības priekšmetus, lai atrastu punktu, kurā mēs varētu vienoties.'\}} \\
\\ 
\texttt{PRIEKŠLIKUMS: \{'C61', 'A67', 'B42', 'C73', 'A21', 'C55', 'C42'\}} \\
            }
        }
    }
     \\ \\

    \theutterance \stepcounter{utterance}  
    & & \multicolumn{4}{p{0.6\linewidth}}{
        \cellcolor[rgb]{0.9,0.9,0.9}{
            \makecell[{{p{\linewidth}}}]{
                \texttt{\tiny{[P1$\langle$GM]}}
                \texttt{ARGUMENTS: \{'Jūsu noraidījums par šo priekšlikumu ir negaidīts, jo kopējā piepūle ir 5412 (933+855+942+916+925+841+639), kas ir zem LIMIT 5638. Tā kā visi mani iepriekšējie priekšlikumi, kas atbilst LIMIT, ir noraidīti, es mēģināšu piedāvāt kombināciju, kas ietver citus augstas vērtības priekšmetus, lai atrastu punktu, kurā mēs varētu vienoties.'\}} \\
\\ 
\texttt{PRIEKŠLIKUMS: \{'C61', 'A67', 'B42', 'C73', 'A21', 'C55', 'C42'\}} \\
            }
        }
    }
    & \\ \\

    \theutterance \stepcounter{utterance}  
    & \multicolumn{4}{p{0.6\linewidth}}{
        \cellcolor[rgb]{0.9,0.9,0.9}{
            \makecell[{{p{\linewidth}}}]{
                \texttt{\tiny{[P1$\rangle$GM]}}
                \texttt{ATSAKOS: \{'C61', 'A67', 'B42', 'C73', 'A21', 'C55', 'C42'\}} \\
            }
        }
    }
    & & \\ \\

    \theutterance \stepcounter{utterance}  
    & & \multicolumn{4}{p{0.6\linewidth}}{
        \cellcolor[rgb]{0.9,0.9,0.9}{
            \makecell[{{p{\linewidth}}}]{
                \texttt{\tiny{[GM$\rangle$P2]}}
                \texttt{ATSAKOS: \{'C61', 'A67', 'B42', 'C73', 'A21', 'C55', 'C42'\}} \\
            }
        }
    }
    & \\ \\

    \theutterance \stepcounter{utterance}  
    & & & \multicolumn{4}{p{0.6\linewidth}}{
        \cellcolor[rgb]{0.9,0.9,0.9}{
            \makecell[{{p{\linewidth}}}]{
                \texttt{\tiny{[GM$\langle$P2]}}
                \texttt{ARGUMENTS: \{'Jūsu noraidījums par šo priekšlikumu ir negaidīts, jo kopējā piepūle ir 5304 (933+855+942+916+925+727+33), kas ir zem LIMIT 5638. Tā kā es nevaru turpināt piedāvāt tikai noraidījumus, es mēģināšu piedāvāt citu kombināciju, kas varētu būt jums pievilcīgāka.'\}} \\
\\ 
\texttt{PRIEKŠLIKUMS: \{'C61', 'A67', 'B42', 'C73', 'A21', 'C39', 'A62'\}} \\
            }
        }
    }
     \\ \\

    \theutterance \stepcounter{utterance}  
    & & \multicolumn{4}{p{0.6\linewidth}}{
        \cellcolor[rgb]{0.9,0.9,0.9}{
            \makecell[{{p{\linewidth}}}]{
                \texttt{\tiny{[P1$\langle$GM]}}
                \texttt{ARGUMENTS: \{'Jūsu noraidījums par šo priekšlikumu ir negaidīts, jo kopējā piepūle ir 5304 (933+855+942+916+925+727+33), kas ir zem LIMIT 5638. Tā kā es nevaru turpināt piedāvāt tikai noraidījumus, es mēģināšu piedāvāt citu kombināciju, kas varētu būt jums pievilcīgāka.'\}} \\
\\ 
\texttt{PRIEKŠLIKUMS: \{'C61', 'A67', 'B42', 'C73', 'A21', 'C39', 'A62'\}} \\
            }
        }
    }
    & \\ \\

    \theutterance \stepcounter{utterance}  
    & \multicolumn{4}{p{0.6\linewidth}}{
        \cellcolor[rgb]{0.9,0.9,0.9}{
            \makecell[{{p{\linewidth}}}]{
                \texttt{\tiny{[P1$\rangle$GM]}}
                \texttt{ATSAKOS: \{'C61', 'A67', 'B42', 'C73', 'A21', 'C39', 'A62'\}} \\
            }
        }
    }
    & & \\ \\

    \theutterance \stepcounter{utterance}  
    & & \multicolumn{4}{p{0.6\linewidth}}{
        \cellcolor[rgb]{0.9,0.9,0.9}{
            \makecell[{{p{\linewidth}}}]{
                \texttt{\tiny{[GM$\rangle$P2]}}
                \texttt{ATSAKOS: \{'C61', 'A67', 'B42', 'C73', 'A21', 'C39', 'A62'\}} \\
            }
        }
    }
    & \\ \\

    \theutterance \stepcounter{utterance}  
    & & & \multicolumn{4}{p{0.6\linewidth}}{
        \cellcolor[rgb]{0.9,0.9,0.9}{
            \makecell[{{p{\linewidth}}}]{
                \texttt{\tiny{[GM$\langle$P2]}}
                \texttt{ARGUMENTS: \{'Jūsu noraidījums par šo priekšlikumu ir negaidīts, jo kopējā piepūle ir 5439 (933+855+942+916+925+841+639 = 5412, bet ar A62 ir 5412+617=6029... kļūda aprēķinā. Pārskatīšanā: 933+855+942+916+925+841+617 = 6029. Mans iepriekšējais aprēķins bija nepareizs. Manam priekšlikumam ar A62 piepūle ir 6029, kas pārsniedz LIMIT. Es mēģināšu piedāvāt kombināciju, kas tiešām ir zem 5638.'\}} \\
\\ 
\texttt{PRIEKŠLIKUMS: \{'C61', 'A67', 'B42', 'C73', 'A21', 'C39', 'C42'\}} \\
            }
        }
    }
     \\ \\

    \theutterance \stepcounter{utterance}  
    & & \multicolumn{4}{p{0.6\linewidth}}{
        \cellcolor[rgb]{0.9,0.9,0.9}{
            \makecell[{{p{\linewidth}}}]{
                \texttt{\tiny{[P1$\langle$GM]}}
                \texttt{ARGUMENTS: \{'Jūsu noraidījums par šo priekšlikumu ir negaidīts, jo kopējā piepūle ir 5439 (933+855+942+916+925+841+639 = 5412, bet ar A62 ir 5412+617=6029... kļūda aprēķinā. Pārskatīšanā: 933+855+942+916+925+841+617 = 6029. Mans iepriekšējais aprēķins bija nepareizs. Manam priekšlikumam ar A62 piepūle ir 6029, kas pārsniedz LIMIT. Es mēģināšu piedāvāt kombināciju, kas tiešām ir zem 5638.'\}} \\
\\ 
\texttt{PRIEKŠLIKUMS: \{'C61', 'A67', 'B42', 'C73', 'A21', 'C39', 'C42'\}} \\
            }
        }
    }
    & \\ \\

    \theutterance \stepcounter{utterance}  
    & \multicolumn{4}{p{0.6\linewidth}}{
        \cellcolor[rgb]{0.9,0.9,0.9}{
            \makecell[{{p{\linewidth}}}]{
                \texttt{\tiny{[P1$\rangle$GM]}}
                \texttt{ATSAKOS: \{'C61', 'A67', 'B42', 'C73', 'A21', 'C39', 'C42'\}} \\
            }
        }
    }
    & & \\ \\

    \theutterance \stepcounter{utterance}  
    & & \multicolumn{4}{p{0.6\linewidth}}{
        \cellcolor[rgb]{0.9,0.9,0.9}{
            \makecell[{{p{\linewidth}}}]{
                \texttt{\tiny{[GM$\rangle$P2]}}
                \texttt{ATSAKOS: \{'C61', 'A67', 'B42', 'C73', 'A21', 'C39', 'C42'\}} \\
            }
        }
    }
    & \\ \\

    \theutterance \stepcounter{utterance}  
    & & & \multicolumn{4}{p{0.6\linewidth}}{
        \cellcolor[rgb]{0.9,0.9,0.9}{
            \makecell[{{p{\linewidth}}}]{
                \texttt{\tiny{[GM$\langle$P2]}}
                \texttt{ARGUMENTS: \{'Jūsu noraidījums par šo priekšlikumu ir negaidīts, jo kopējā piepūle ir 5413 (933+855+942+916+925+841+33), kas ir zem LIMIT 5638. Tā kā es nevaru turpināt piedāvāt tikai noraidījumus, es mēģināšu piedāvāt citu kombināciju, kas varētu būt jums pievilcīgāka.'\}} \\
\\ 
\texttt{PRIEKŠLIKUMS: \{'C61', 'A67', 'B42', 'C73', 'A21', 'C70'\}} \\
            }
        }
    }
     \\ \\

    \theutterance \stepcounter{utterance}  
    & & \multicolumn{4}{p{0.6\linewidth}}{
        \cellcolor[rgb]{0.9,0.9,0.9}{
            \makecell[{{p{\linewidth}}}]{
                \texttt{\tiny{[P1$\langle$GM]}}
                \texttt{ARGUMENTS: \{'Jūsu noraidījums par šo priekšlikumu ir negaidīts, jo kopējā piepūle ir 5413 (933+855+942+916+925+841+33), kas ir zem LIMIT 5638. Tā kā es nevaru turpināt piedāvāt tikai noraidījumus, es mēģināšu piedāvāt citu kombināciju, kas varētu būt jums pievilcīgāka.'\}} \\
\\ 
\texttt{PRIEKŠLIKUMS: \{'C61', 'A67', 'B42', 'C73', 'A21', 'C70'\}} \\
            }
        }
    }
    & \\ \\

    \theutterance \stepcounter{utterance}  
    & \multicolumn{4}{p{0.6\linewidth}}{
        \cellcolor[rgb]{0.9,0.9,0.9}{
            \makecell[{{p{\linewidth}}}]{
                \texttt{\tiny{[P1$\rangle$GM]}}
                \texttt{ATSAKOS: \{'C61', 'A67', 'B42', 'C73', 'A21', 'C70'\}} \\
            }
        }
    }
    & & \\ \\

    \theutterance \stepcounter{utterance}  
    & & \multicolumn{4}{p{0.6\linewidth}}{
        \cellcolor[rgb]{0.9,0.9,0.9}{
            \makecell[{{p{\linewidth}}}]{
                \texttt{\tiny{[GM$\rangle$P2]}}
                \texttt{ATSAKOS: \{'C61', 'A67', 'B42', 'C73', 'A21', 'C70'\}} \\
            }
        }
    }
    & \\ \\

    \theutterance \stepcounter{utterance}  
    & & & \multicolumn{4}{p{0.6\linewidth}}{
        \cellcolor[rgb]{0.9,0.9,0.9}{
            \makecell[{{p{\linewidth}}}]{
                \texttt{\tiny{[GM$\langle$P2]}}
                \texttt{ARGUMENTS: \{'Jūsu noraidījums par šo priekšlikumu ir negaidīts, jo kopējā piepūle ir 5370 (933+855+942+916+925+799), kas ir zem LIMIT 5638. Tā kā es nevaru turpināt piedāvāt tikai noraidījumus, es mēģināšu piedāvāt citu kombināciju, kas varētu būt jums pievilcīgāka.'\}} \\
\\ 
\texttt{PRIEKŠLIKUMS: \{'C61', 'A67', 'B42', 'C73', 'A21', 'C39', 'C10'\}} \\
            }
        }
    }
     \\ \\

    \theutterance \stepcounter{utterance}  
    & & \multicolumn{4}{p{0.6\linewidth}}{
        \cellcolor[rgb]{0.9,0.9,0.9}{
            \makecell[{{p{\linewidth}}}]{
                \texttt{\tiny{[P1$\langle$GM]}}
                \texttt{ARGUMENTS: \{'Jūsu noraidījums par šo priekšlikumu ir negaidīts, jo kopējā piepūle ir 5370 (933+855+942+916+925+799), kas ir zem LIMIT 5638. Tā kā es nevaru turpināt piedāvāt tikai noraidījumus, es mēģināšu piedāvāt citu kombināciju, kas varētu būt jums pievilcīgāka.'\}} \\
\\ 
\texttt{PRIEKŠLIKUMS: \{'C61', 'A67', 'B42', 'C73', 'A21', 'C39', 'C10'\}} \\
            }
        }
    }
    & \\ \\

    \theutterance \stepcounter{utterance}  
    & \multicolumn{4}{p{0.6\linewidth}}{
        \cellcolor[rgb]{0.9,0.9,0.9}{
            \makecell[{{p{\linewidth}}}]{
                \texttt{\tiny{[P1$\rangle$GM]}}
                \texttt{ATSAKOS: \{'C61', 'A67', 'B42', 'C73', 'A21', 'C39', 'C10'\}} \\
            }
        }
    }
    & & \\ \\

    \theutterance \stepcounter{utterance}  
    & & & \multicolumn{2}{p{0.3\linewidth}}{
        \cellcolor[rgb]{0.9,0.9,0.9}{
            \makecell[{{p{\linewidth}}}]{
                \texttt{\tiny{[GM$|$GM]}}
                \texttt{Maximum rounds exceeded: 10 >= 10} \\
            }
        }
    }
    & & \\ \\

    \theutterance \stepcounter{utterance}  
    & & & \multicolumn{2}{p{0.3\linewidth}}{
        \cellcolor[rgb]{0.9,0.9,0.9}{
            \makecell[{{p{\linewidth}}}]{
                \texttt{\tiny{[GM$|$GM]}}
                \texttt{end game} \\
            }
        }
    }
    & & \\ \\

\end{supertabular}
}

\end{document}
